
\documentclass[screen]{memarticle}
%\documentclass[draft,draftwater]{memarticle}
%\documentclass{memarticle}

%
% Add other packages here, as needed.
%
\usepackage{tikz-uml}
\tikzumlset{fill class=yellow!10}

%------------------------------------------------------------
% Document title control. Change these as appropriate for your
% document.
%------------------------------------------------------------
\newcommand{\doctitle}{Polymorphism and Data Products}
\newcommand{\brieftitle}{\doctitle}
\newcommand{\authors}{\href{mailto:paterno@fnal.gov}{Marc Paterno}}
\newcommand{\docversion}{1}

%------------------------------------------------------------
% Define new commands here
%------------------------------------------------------------

\newcommand{\art}{\name{art}\xspace}
\lstMakeShortInline[language=isocpp,style=cpplistingstyle]|

%------------------------------------------------------------
% Document starts here
%------------------------------------------------------------
\begin{document}

\topmatter
	
\screenflush\begin{abstract}
  This document explains why the \art development team discourages the
  use of base class pointers in data products. It also explains the
  rationale for, and the use of, ``views'', as provided by the \art
  framework.
\end{abstract}

\chapter{Introduction\label{ch:introduction}}

In \cpp, \emph{polymorphism}, or, to be more precise \emph{run-time
  polymorphism},\footnote{Run-time polymorphism is to be distinguished
  from \emph{parametric polymorphism}, also called \emph{compile-time
    polymorphism}, which in \cpp is provided by class and function
  templates.} is the mechanism, using |virtual| functions, through
which a base class can declare a function and various derived classes
can differently implement that function. The actual function chosen be
executed is determined by the run-time type of the object on which the
function call is made.

\emph{Type erasure} is the compile-time process through which a
program looses some (or all) type information about a given object. In
\cpp, a program that uses run-time polymorphism frequently also uses
type erasure---specifically, it does so whenever the concrete type of
an object has any function members not in the interface of that
concrete types base class. Any such functions become unavailable to
code that knows only about the base class, unless a |dynamic_cast| is
used to recover the hidden type information.

In \art, recommended practice is the following:
\begin{enumerate}
  \item\label{no-erasure} data products should avoid type erasure, and
  \item\label{poly-algor} polymorphic use of data products should be achieved through
    algorithms that manipulate base class pointers or references.
\end{enumerate}
Item~\ref{no-erasure} means that data products should only contain
concrete types, preferably contained by value (not by reference); they
should \emph{not} contain pointers-to-base, especially pointers which
are used to implement ownership. Item~\ref{poly-algor} means that many
algorithms should be written as functions which accept
pointers-to-base, or collections of pointers-to-base (with appropriate
|const| qualification). We supply the function templates
|Event::getView|, and the class template |View<E>| to support this
use.

This document explains the rationale for, and the behavior of,
|Event::getView| and |View<E>|, and recommended practice on the design
of data products and the design of algorithms that are polymorphic on
those data products.

\chapter{The problem with type erasure}

A container such as |std::vector<Base*>| is unsuitable for use as a
data product in \art because it (the container) does not \emph{own}
the objects to which the elements point. This means that when the
|vector| is deleted the pointed-to objects are \emph{not}
deleted---and thus are leaked. For proper functioning, \art assumes
that data products own their contents.

One attempt to ``fix'' this problem---which we do not endorse---is to
create a class template which, like |vector<Base*>|, is capable of
holding any class that inherits from |Base|, but which, unlike
|vector<Base*>|, owns the objects to which it points. This elimantes
the problem of memory leaks, but it leads to other problems. These
include:
\begin{itemize}
\item frequent use of |dynamic_cast|,
\item complexity of management of ownership of the pointed-to elements, and
\item complexity induced in |Ptr<E>|.
\end{itemize}

\section{Frequent use of \lstinline{dynamic_cast}}

In code that uses type-erasing containers, it is common to find
widespread use of |dynamic_cast|. This is because it is very rare for
\emph{all} client code that uses the type-erased contained to be
independent of the derived classes. If


\section{Management of ownership}
\section{Complexity in \lstinline{Ptr<T>}}

\chapter{What is a ``view''?}

In \art, we use the term ``view'' to denote \emph{a collection of
  pointers to objects contained by collections, which collections are
  data products in an \lstinline{Event}}.  \art supports two different
class templates as ``views'': |std::vector<E const*>| and
|art::View<E>|.

The most efficient type, for most purposes, %
is |std::vector<const T*>|. Note that the elements in the |vector| are
pointers-to-const. The objects being pointed to (the objects of type
|T|, which are elements in some collection in the event) may not be
modified through these pointers; if the objects are of class-type,
only |const| member functions can be called on them. The pointers may
be modified (\textit{e.g.}, a pointer may be overwritten by a
different value such as null), or the collection may be modified
(\textit{e.g.}, it may be sorted, or elements can be removed), or
both. The type |T| may be to the base class of the actual, concrete
class of the objects stored in the |Event|; see
\S\ref{ch:concrete-use} for an example. Thus the collection of
pointers may be used polymorphically by an algorithm that depends only
on the base class |T|, not on the (unspecified) exact type of the
object being pointed to. It is important to note that the view
|vector| \emph{does not own} the objects to which it points. These
objects are not deleted when the |vector| is destroyed. One should not
call |delete| on the contained pointers, or use one as the value
passed to any function or constructor that takes ownership of the
pointed-to object, \textit{e.g.}, the constructor of
|std::shared_ptr|. Because it does not own the objects to which it
points, such a |vector| is not suitable for persistence, and should
2not be put into an |Event|.

The class template |art::View<E>| should be used when one needs the
ability to create a persistent object from a ``view''. At the cost of
some additional overhead (when compared to |std::vector<E const*>|),
|View<E>| provides the ability to create a |art::PtrVector<E>|. Like
|std::vector<E const*>|, |View<E>| provides only |const| access to the
objects pointed to by the ``view''. See \S\ref{ch:view-template} for
more details on |art::View<E>|.

\chapter{Obtaining ``views''}

The class |Event| has an overload set\footnote{In \cpp, an
  \emph{overload set} is a set of functions, function templates, or
  both, with the same \emph{name}, but with different
  \emph{types}. The type of a function is determined by its list of
  parameters and its return type; the type of a function template is
  determined by its list of function parameters, the list of template
  parameters, and its return type. Functions and function templates in
  an overload set are not permitted to differ only in return type.}
named |getView|. There are four functions in the overload set; they
differ in how they specify the collection into which the view is
pointing (with two options: either a combination of module label and
product instance name, or an |InputTag|) and by the class template
used to contain the view (with two choices, either %
|std::vector<E const*>| or |View<E>|). These function templates are shown in
listing~\ref{l:getview}.
\begin{lstlisting}[language=C++%
                   , caption={The \lstinline{Event::getView} member function templates.}%
                   , label={l:getview}
                   , float=htb
                   ]
template <class E>
std::size_t Event::getView(std::string const& moduleLabel,
                           std::string const& productInstanceName,
                           std::vector<E const*>& v) const;
template <class E>
std::size_t Event::getView(InputTag const& t, 
                           std::vector<E const*>& v) const;
template <class E>
bool Event::getView(std::string const& moduleLabel,
                    std::string const& productInstanceName,
                    View<E>& v) const;
template <class E>
bool Event::getView(InputTag const& t, View<E>& v) const;
\end{lstlisting}

For each member function template

\chapter{The class template \lstinline{View<E>}\label{ch:view-template}}

Listing~\ref{l:view} shows the main features of the class template |art::View<E>|.

\begin{lstlisting}[language=C++%
                   , caption={The public interface of class template \lstinline{art::View<E>}.}%
                   , label={l:view}
                   , float=htb
                   ]
template <class E> class art::View {
public:
  typedef std::vector<E const*> collection_type;
  View(); // default constructed Views are invalid
  bool isValid() const;
  collection_type&       vals();
  collection_type const& vals() const;
  ProductID id() const;
  void fill(PtrVector<T>& pv) const;

  // The following are conversion operators to collection_type& and
  // collection_type const&. These allow an object of type View<E> to
  // be passed to any function that expects a std::vector<E const*>&
  // or std::vector<E const*> const&.
  operator collection_type& ();
  operator collection_type const& () const;
};
\end{lstlisting}

\chapter{A concrete use\label{ch:concrete-use}}
The UML class diagram in figure~\ref{fig:prongdiagram} shows the
relationship between the classes |Trajectory|, |MCParticle|, and
|Track|. In our example, |Trajectory| is an \emph{abstract}
class. This means it contains at least one \emph{pure} virtual
function---a virtual function with no implementation. Because it is
abstract, no instances of class |Trajectory| can be created. Instead,
|Trajectory| is useful as a \emph{base class}, because it provides an
interface that can be used polymorphically---by lstinline that depends on
the base class |Trajectory|, but does not depend on knowing about any
derived types. |Trajectory| contains one interesting member function,
the virtual function |momentum(double s)|, which given a position |s|
along the length of the trajectory, will provide the momentum of the
represented object at that location. |Trajectory| has two subclasses,
also called derived classes: |MCParticle| (which represents a Monte
Carlo particle) and |Track| (which represents a reconstructed track in
a detector).
\begin{figure}[!ht]
  \centering
  \begin{tikzpicture}
    \umlclass[type=abstract]{Trajectory}{}{\umlvirt{momentum(double s) : ThreeVector} \\ \umlvirt{$\sim$Trajectory()} \\ \ldots}
    \umlclass[x=-4, y=-4]{MCParticle}{}{momentum(double) : ThreeVector \\ $\sim$MCParticle() \\ \ldots}
    \umlclass[x= 4, y=-4]{Track}{}{momentum(double s) : ThreeVector \\ $\sim$Track() \\ \ldots}

    \umlinherit[geometry=|-|]{MCParticle}{Trajectory}
    \umlinherit[geometry=|-|]{Track}{Trajectory}
  \end{tikzpicture}
  \caption{The UML class diagram for our example problem. The lack of
    an argument name in the function of \lstinline{MCParticle} is meant to
    convey the information that this function does not use the value of the argument.}
  \label{fig:prongdiagram}
\end{figure}

\appendix

% \chapter{References\label{ch:references}}

% Include as needed.

% \chapter{Appendices\label{ch:appendices}}

% Include as needed.

% Include '\bibname', and the bibliography parts below, only if you
% have a bibliography. Not every tech report needs to have one.

%\bibname

% The 'thebibliography' environment takes one required argument, which
% is a piece of text as wide as the widest label in the
% bibliography. Ugly, but true.
%\begin{thebibliography}{9}

% \bibitem{cpp:2003}
%  \textit{ISO/IEC 14882:2003 (Second Edition) Programming Languages - C++},
%   John Wiley \& Sons, Ltd.

%\end{thebibliography}

\end{document}
